\documentclass[11pt,a4paper]{article}
\usepackage[margin=22mm,headheight=14pt]{geometry}
\usepackage[T1]{fontenc}\usepackage[utf8]{inputenc}
\usepackage{lmodern,microtype,amsmath,amssymb,bm,booktabs,longtable,array,hyperref,fancyhdr}
\hypersetup{colorlinks=true,linkcolor=black,urlcolor=black}
\pagestyle{fancy}\fancyhf{}\lhead{QICT: Predictions and Derived Results}\rhead{Domains and supporting tests}\cfoot{\thepage}
\setlength{\parskip}{5pt}\setlength{\parindent}{13pt}\setlength{\emergencystretch}{5em}
\newcommand{\dd}{\mathrm{d}}\newcommand{\Mp}{\,\mathrm{Mpc}}\newcommand{\msun}{M_\odot}
\newcommand{\Gev}{\,\mathrm{GeV}}\newcommand{\Om}{\Omega}\newcommand{\sigv}{\langle\sigma v\rangle}
\title{\bfseries Theoretical Predictions and Derived Results\\[6pt]\large Spectral Tests of Dark Matter and Cosmic Structure in Quantum Information Copy Time}
\author{Sacha Mohamed\\\small Casablanca, Morocco}\date{19 October 2026}
\begin{document}\maketitle

This register collects the numbered predictions, derived relations and numerical benchmarks of the main article and the supplementary identities. Each entry specifies its domain and supporting calculation. Cross-references give the section, proposition, equation or table number in the compiled document named in the same row.

The entries distinguish finite QICT operator certificates, conditional physical identities, declared model parameters and calculations using observational reference inputs. The scalar-portal masses and couplings, background normalization and halo profiles retain their stated external provenance. A numerical result for a declared reference model has that model as its domain. The finite C9 certificate concerns the static chiral-source polar-unitary construction described in the main article.

The supplied finite calculations establish the spectral and response results listed below. Their physical extension uses the source, state and convergence hypotheses in Supplementary Sections S1--S4. In particular, a source-visible pole requires a nonzero residue in addition to an isolated operator eigenvalue; SI10 and SI11 give complementary linear and amplitude enclosures under verified common-limit error bounds. A QICT dark-particle mass, its relic abundance, complete primordial CMB spectra and a sample-wide galaxy fit are not established by these inputs. The support column records the tests actually performed; mathematical or numerical agreement is distinct from independent experimental confirmation.

\small
\setlength{\tabcolsep}{3pt}\renewcommand{\arraystretch}{1.15}
% Generated from the main register and verified supplementary identities.
\newcommand{\predictionrows}{%
\multicolumn{4}{@{}l@{}}{\textbf{Source spectra and finite operator certificates}}\\*
SP1 & $\lambda_{\max}(F_{13}(0))=0.016084$ at R4 & Main article: Table~1; Six-source C9 Fej\'er upper bound; finite regulator & Finite moments; exact recorded-data certificate\\
SP2 & $\lambda_{\max}(B_{13}(0))\simeq0.001956273$ at R4 & Main article: Table~2; Inverse-Toeplitz upper bound on the same putative atom & Finite moments; exact recorded-data certificate\\
SP3 & R4 atomic weight $\leq0.001956273015$ & Main article: Table~3; Exact rational, normalized recorded C9 moments & Finite moments; exact recorded-data certificate\\
SP4 & Flat moment rank reconstructs all cyclic dynamics & Main article: Proposition~10; exact positivity and rank equality & Analytic proof; numerical matrix controls\\
SP5 & $\|C_n-C_1^n\|\leq\min\{2,n(n-1)\varepsilon/2\}$ & Main article: Proposition~11; unitary transfer and isometric source & Analytic proof; numerical matrix controls\\
SP6 & Exponential atomic-weight enclosure & Main article: Proposition~14; independently certified angular gap & Analytic proof; numerical matrix controls\\
SP7 & Joint regulator and moment limit & Main article: Section~2.13; uniform gap and weak convergence & Analytic proof under the stated hypotheses\\
SP8 & Finite-arc spectral-weight bound & Main article: Eq.~(2.17); stated angular resolution & Analytic proof under the stated hypotheses\\
SP9 & Fixed-polynomial moment-error enclosure & Main article: Eq.~(2.18); verified moment errors & Analytic proof under the stated hypotheses\\
SP10 & Optimal bound decreases to the visible atom & Main article: Proposition~13; fixed positive measure & Analytic proof under the stated hypotheses\\
SP11 & Largest feasible atomic component is $B_N$ & Main article: Proposition~12; positive-definite Toeplitz form & Analytic proof under the stated hypotheses\\
SP12 & $\mu_L(\{1\})=\mu_R(\{1\})=0$ at R1--R4 & Main article: Proposition~17; certified polar-unitary finite model & Analytic proof; outward-rounded operator certificate\\
SP13 & R4 chord gaps $g_L\geq1.92308\times10^{-7}$, $g_R\geq3.19912\times10^{-4}$ & Main article: Table~4; all relevant spectral pairs & Analytic proof; outward-rounded operator certificate\\
SP14 & $F_N(0)\leq\min\{1,4/(N^2g^2)\}$, $|G(1)|\leq g^{-1}$ & Main article: Eq.~(3.8); normalized separated finite measure & Analytic proof; outward-rounded operator certificate\\
SP15 & Operator errors with $\sum_j\eta_j<g$ preserve separation & Main article: Section~3.3; fixed source and spectator structure & Analytic proof; outward-rounded operator certificate\\
SP16 & Positive finite Lehmann weights reconstruct every moment & Main article: Proposition~15; stationary state and isometric source & Analytic proof; independent moment reconstruction\\
\multicolumn{4}{@{}l@{}}{\textbf{Identifiability and independent operator information}}\\*
MAT1 & Distinct spectral types with identical moments through any finite order $N$ & Main article: Proposition~1; explicit normalized scalar and six-channel measures & Analytic proof or explicit construction\\
MAT2 & Integer energy branches share one transfer unitary & Main article: Proposition~5; fixed clock & Analytic proof or explicit construction\\
MAT3 & Identical two-point moments admit distinct four-point interactions & Main article: Proposition~6; explicit positive Hamiltonian family & Analytic proof or explicit construction\\
MAT4 & Invisible odd-sector phases are independently variable & Main article: Proposition~3; sources supported in the even sector & Analytic proof or explicit construction\\
MAT5 & Clock scaling gives $E\mapsto E/s$ & Main article: Proposition~2; fixed dimensionless transfer & Analytic proof or explicit construction\\
MAT6 & $\|\rho_{\mathcal G}-\rho_{\rm diag}\|_1\leq\varepsilon_{\rm coh}$ & Main article: Eq.~(3.3); retained joint eigenspace coherences & Analytic proof or explicit construction\\
\multicolumn{4}{@{}l@{}}{\textbf{Finite gravitational response}}\\*
GR1 & $\epsilon_{0,\mathrm{src},R}=5.889\times10^{-5}$ & Main article: Section~4.4; Finite source convention: $\partial_J\ln\sqrt{-g}=4$ & Nine-state source evaluation\\
GR2 & Connected $D(0)=0.1494933611$; curvature $0.05084820378$ & Main article: Equation~(4.8); Nine-state source derivatives, distinct contact term & Nine-state source evaluation\\
GR3 & Positive response derivatives and alternating remainder & Main article: Proposition~20; positive Lehmann weights & Analytic proof; finite-response checks\\
GR4 & $a_1=0.35360648$, $a_2=1.29642815$, analytic disk $|s|<0.2209$ & Main article: Eq.~(4.11); supplied nine-state source & Nine-state source evaluation\\
GR5 & Contact shift $B\mapsto B+bI$ shifts curvature by $b$ & Main article: Eq.~(4.13); fixed $K_0$ and linear source & Analytic proof; finite-response checks\\
GR6 & Curvature error $\leq Mx+(\delta_h+\delta_a)/x$ & Main article: Eq.~(4.15); declared remainder and data errors & Analytic proof; finite-response checks\\
\multicolumn{4}{@{}l@{}}{\textbf{Gravitational matching and expansion}}\\*
DE1 & $G=c^5\tau_*^2/(8\pi\hbar\gamma_R)$ & Main article: Equation~(4.2); Conditional QICT identity & Conditional physical identity\\
DE2 & $\Lambda=\epsilon_0/(\gamma_Rc^2\tau_*^2)$ & Main article: Equation~(4.2); Conditional QICT identity & Conditional physical identity\\
DE3 & $\Lambda_{\rm obs}=1.09076\times10^{-52}$ m$^{-2}$ & Main article: Equation~(6.2); Observational-input reconstruction & Reconstruction from stated observational inputs\\
DE4 & $u_{\Lambda,\rm obs}=5.2525\times10^{-10}$ J/m$^3$ & Main article: Equation~(6.2); Observational-input reconstruction & Reconstruction from stated observational inputs\\
DE5 & $p=-u_{\rm vac}$; $w=-1$ for $u_{\rm vac}\ne0$ & Main article: Proposition~19; Conditional metric-response theorem & Conditional physical identity\\
DE6 & $\epsilon_0/\gamma_R^2=(7.1528\pm0.1910)\times10^{-121}$ & Main article: Equation~(4.5); Planck base-$\Lambda$CDM reference & Reconstruction from stated observational inputs\\
\multicolumn{4}{@{}l@{}}{\textbf{Particle interactions and thermal production}}\\*
DM1 & $m_S=10,40,300$ GeV (distinct examples) & Main article: Section~5; Independent effective-model inputs & Declared effective-model input\\
DM2 & $\sigma_{\rm SI}(40)=5.282\times10^{-46}$ cm$^2$ & Main article: Section~5.1; Higgs portal benchmark; excluded; full local scalar abundance & Rate calculation; published LZ comparison\\
DM3 & $\sigma_{\rm SI}(300)=3.910\times10^{-45}$ cm$^2$ & Main article: Section~5.1; Higgs portal benchmark & Specified thermal or scattering model calculation\\
DM4 & $\sigv_{hh}^{\rm contact}=2.345\times10^{-26}$ cm$^3$/s & Main article: Section~5.2; Higgs-pair contact contribution & Specified thermal or scattering model calculation\\
DM5 & $Y_{\rm FI}=3.6292\times10^{-11}$ & Main article: Equation~(5.8); One daughter per parent decay; stated parent, width and thermal history & Specified thermal or scattering model calculation\\
DM6 & $\Om_{\rm FI}h^2=0.09951$ & Main article: Section~5.3; Freeze-in, selected effective width & Specified thermal or scattering model calculation\\
DM7 & $|\lambda_{HS}|\leq6.454\times10^{-4}$ at 40 GeV & Main article: Equation~(5.11); Inversion of published LZ bound in selected model & Rate calculation; published LZ comparison\\
DM8 & $\Gamma_A^{\rm req.}=1.205877\times10^{-20}$ GeV & Main article: Equation~(5.12); Observation-normalized inverse calculation & Specified thermal or scattering model calculation\\
DM9 & Constant-rate freeze-out $\Omega_\chi h^2=0.100322$ & Main article: Equation~(5.15); Numerical Boltzmann limiting equation & Specified thermal or scattering model calculation\\
DM10 & $Y_\infty\propto\delta_A\sqrt{8\pi\gamma_R}/\varepsilon_A^2$ & Main article: Equation~(5.28); Exact scale cancellation in specified one-daughter limit & Analytic derivation; thermal numerical controls\\
DM11 & $\Omega_Sh^2=0.0234522$ at $m_S=300$ GeV, $\lambda_{HS}=0.20$ & Main article: Equation~(5.22); Thermal tree-level two-body portal, constant $g_*=90$ & Specified thermal or scattering model calculation\\
DM12 & $I_0^{-1}\leq Y_\infty\leq(I_0-E)^{-1}$ & Main article: Proposition~21; positive source and $E<I_0$ & Analytic derivation; thermal numerical controls\\
DM13 & Decay yield with variable thermal degrees of freedom & Main article: Eq.~(5.32); entropy conservation and equilibrated parent & Analytic derivation; thermal numerical controls\\
DM14 & $\partial\ln Y/\partial\ln T_{\rm RH}=7.78357$ at $x_{\rm RH}=10$ & Main article: Eq.~(5.33); constant-degree decay model & Analytic derivation; thermal numerical controls\\
DM15 & Odd-branch stability under a strict final-state threshold & Main article: Proposition~4; exact parity and energy conservation & Analytic stability proof\\
\multicolumn{4}{@{}l@{}}{\textbf{Matter growth and structure}}\\*
LSS1 & $D(0.5)=0.76894$ & Main article: Section~6.2.1; Smooth-$\Lambda$ baseline & Reference-model numerical calculation\\
LSS2 & $f\sigma_8(0.5)=0.47449$ & Main article: Section~6.2.1; Baseline, observed $\sigma_8$ input & Reference-model numerical calculation\\
LSS3 & $\dd n/\dd\ln M(10^{14}\msun)=2.704\times10^{-5}$ Mpc$^{-3}$ & Main article: Section~6.3.1; Approximate Press--Schechter & Reference-model numerical calculation\\
LSS4 & Spherical collapse threshold $a=0.4614218$ & Main article: Section~6.3.2; Chosen overdensity and mass & Reference-model numerical calculation\\
LSS5 & $C_{50}^{\kappa\kappa}=1.701\times10^{-8}$; $C_{200}^{\kappa\kappa}=2.544\times10^{-9}$ & Main article: Equation~(7.12); Single-source Limber/linear GR baseline & Reference-model numerical calculation\\
\multicolumn{4}{@{}l@{}}{\textbf{Galaxy populations and gas dynamics}}\\*
GAL1 & $M_\star(10^{12}\msun)=7.816\times10^9\msun$ & Main article: Section~6.3.3; Chosen efficiency $\varepsilon_\star$ & Reference-model numerical calculation\\
GAL2 & $\xi_g(10\Mp)=1.462405$ for $b_g=1.5$ & Main article: Section~6.3.3; Smooth-spectrum benchmark & Reference-model numerical calculation\\
GAL3 & $M_g(3\,\mathrm{Gyr})=8.8841\times10^9M_\odot$ & Main article: Equation~(6.19); Selected one-zone gas-regulator inputs & Reference-model numerical calculation\\
GAL4 & $M_*(3\,\mathrm{Gyr})=5.4619\times10^9M_\odot$ & Main article: Equation~(6.19); Same gas-regulator solution & Reference-model numerical calculation\\
GAL5 & Positive gas solution for time-dependent inflow & Main article: Eq.~(6.20); nonnegative rates and stated return fraction & Exact reservoir solution and balance check\\
GAL6 & $M_w(3\,\mathrm{Gyr})=9.10314\times10^9M_\odot$ & Main article: Eq.~(6.21); exact baryon balance in the selected reservoir & Exact reservoir solution and balance check\\
\multicolumn{4}{@{}l@{}}{\textbf{Particle--mesh evolution}}\\*
PM1 & Final density RMS $10.4690$; peak $1484.30$ & Main article: Section~6.4; Coherent-wave periodic $32^3$ PM experiment & 200/300-step comparison at the stated mesh\\
\multicolumn{4}{@{}l@{}}{\textbf{Lensing and scalar metric response}}\\*
LEN1 & SIS $\theta_E=0.8884$ arcsec & Main article: Section~7.2; Chosen geometry and dispersion & Reference-model numerical calculation\\
LEN2 & SIS $|\Delta t|=42.93$ d & Main article: Section~7.2; Chosen source offset & Reference-model numerical calculation\\
LEN3 & NFW $\kappa(0.1\Mp)=0.1197$ & Main article: Section~7.3; Chosen cluster mass/concentration & Reference-model numerical calculation\\
LEN4 & NFW $\gamma_t(0.1\Mp)=0.0708$ & Main article: Section~7.3; Chosen cluster mass/concentration & Reference-model numerical calculation\\
LEN5 & Cluster $\theta_t=5.35$ arcsec & Main article: Section~7.3; Separate idealized massive cluster & Reference-model numerical calculation\\
LEN6 & Angular convergence ratio at $\ell=200$: $1.2709505$ & Main article: Section~7.5; Selected scalar response and matched initial density & Reference-model numerical calculation\\
LEN7 & $E_G(0.5)=0.41400$ (GR), $0.40097$ (selected scalar response) & Main article: Section~7.5; Declared background and growth conditions & Reference-model numerical calculation\\
LEN8 & Growth and $E_G$ invert to $\mu$ and $\eta$ & Main article: Eq.~(7.17); quasistatic conserved matter and fixed background & Analytic inversion; integrated-growth check\\
\multicolumn{4}{@{}l@{}}{\textbf{CMB geometry and transfer relations}}\\*
CMB1 & $r_s(z_*)=144.294$ Mpc & Main article: Section~8.1; Reference recombination input & Reference-background numerical calculation\\
CMB2 & $100\theta_*=1.04075$ & Main article: Section~8.1; Reference, observational normalization & Reference-background numerical calculation\\
CMB3 & $\ell_A=301.86$ & Main article: Section~8.1; Acoustic spacing of the stated background & Reference-background numerical calculation\\
CMB4 & Saha $z_{x_e=1/2}=1369$ & Main article: Section~8.2; Equilibrium-only ionization & Reference-background numerical calculation\\
CMB5 & $C_\ell^{TT},C_\ell^{TE},C_\ell^{EE},C_L^{\phi\phi}$ & Main article: Equations~(8.9), (8.13); Analytic observables; primordial and transfer inputs specified separately & Transfer-functional identity or analytic proof\\
CMB6 & $|C_\ell^{TE}|^2\leq C_\ell^{TT}C_\ell^{EE}$ & Main article: Proposition~22; positive primordial measure & Transfer-functional identity or analytic proof\\
CMB7 & Finite spectra admit a primordial equivalence class & Main article: Eq.~(8.15); fixed transfers and positive reference spectrum & Transfer-functional identity or analytic proof\\
\multicolumn{4}{@{}l@{}}{\textbf{Supplementary physical identities}}\\*
SI1 & $\|P_{\mathcal R}-P\|\leq r\eta_{\mathcal R}$; isolated rank-one eigenvalue persists & Supplement: Proposition~S1; common physical resolvents; source pole requires positive residue & Contour proof; matrix-sequence controls\\
SI2 & $|E_{\chi,\mathcal R}-E_\chi|\leq2r^2\eta_{\mathcal R}$; $\|Z_{\mathcal R}-Z\|\leq r\eta_{\mathcal R}+2\varepsilon_{\mathcal R}$ & Supplement: Equations~(S2.3)--(S2.4); positive limiting source weight & Analytic bounds; independent contour integration\\
SI3 & Connected cumulants determine $\Gamma^{(3)}$ and $\Gamma^{(4)}$ & Supplement: Proposition~S2; locally invertible source-to-field map & Exact cumulants; independent inverse-map derivatives\\
SI4 & Fixed dynamics admit distinct initial abundances & Supplement: Proposition~S3; specified admissible initial densities & Convex density-operator construction\\
SI5 & $h^2+k\ell_*^2/a^2=(\widehat u+\epsilon_0)/(3\gamma_R)$ & Supplement: Equation~(S4.1); matched Einstein branch and cosmological state & Algebraic substitution of gravitational matching\\
SI6 & $v_c^2=r\Psi^{\prime}$; deflection from $\Phi+\Psi$ & Supplement: Equation~(S4.3); common weak-field metric and stated symmetries & Timelike and null-geodesic relations\\
SI7 & On-shell amplitude from the connected four-point function and pole residues & Supplement: Equation~(S3.5); stable scalar branch, Lorentzian continuation and asymptotic states & LSZ reduction with the specified field normalization\\
SI8 & $\mathrm d\sigma=\overline{|\mathcal M|^2}\mathrm d\Phi_f/(4F S_f)$; $F=\sqrt{(p_1\cdot p_2)^2-m_1^2m_2^2}$ & Supplement: Equation~(S3.6); specified incoming particles and final-state symmetry factor & Relativistic flux and phase-space identity\\
SI9 & Primordial curvature covariance determines the power-spectrum input to CMB transfers & Supplement: Equation~(S4.4); homogeneous isotropic initial state and specified transfer evolution & Spectral covariance definition and transfer integrals\\
SI10 & $u^\dagger Zu\geq u^\dagger Z_{\mathcal R}u-r\eta_{\mathcal R}-2\varepsilon_{\mathcal R}$ & Supplement: Equation~(S2.5); unit source vector and verified common-limit error bounds & Analytic enclosure; visible and invisible source controls\\
SI11 & $\begin{gathered}[\max\{0,\sqrt{z_{\mathcal R}}-d_{\mathcal R}\}]^2\leq z\\z\leq\min\{1,(\sqrt{z_{\mathcal R}}+d_{\mathcal R})^2\}\end{gathered}$ & Supplement: Equation~(S2.6); normalized source; $z=u^\dagger Zu$, $d_{\mathcal R}=r\eta_{\mathcal R}+\varepsilon_{\mathcal R}$ & Reverse-triangle proof; weak-overlap matrix control\\
}

\begin{longtable}{@{}>{\raggedright\arraybackslash}p{.065\linewidth}>{\raggedright\arraybackslash}p{.365\linewidth}>{\raggedright\arraybackslash}p{.335\linewidth}>{\raggedright\arraybackslash}p{.165\linewidth}@{}}
\toprule ID & Statement or value & Domain and cross-reference & Support\\\midrule\endfirsthead
\toprule ID & Statement or value & Domain and cross-reference & Support\\\midrule\endhead
\bottomrule\endlastfoot
\predictionrows
\end{longtable}
\normalsize

\paragraph{Reproducibility and error conventions.}
The numerical values, finite matrix arrays, interval certificates and formula checks are distributed in the code and data folder. Its \texttt{README.md} gives the reproduction commands and expected outputs. The exact recorded-moment certificate, outward-rounded finite-operator enclosure and floating-point numerical comparisons retain their respective arithmetic guarantees. The conditional physical identities carry the assumptions stated in their proofs. Regulator convergence and empirical error bounds enter where explicitly specified; the finite tests retain their stated regulator and model domains.

\end{document}
